\chapter{Laser Photocoagulation of the Retina}

\section*{Overview}
Laser photocoagulation of the retina is an ophthalmic procedure in which a laser is used to seal leaking blood vessels, treat new abnormal vessels, or create burns in the retina to prevent vision loss in conditions such as diabetic retinopathy and retinal tears.
\section*{Alternative Names}
Retinal laser photocoagulation; laser therapy for the retina; panretinal photocoagulation; focal laser treatment
\section*{Principle}
Laser energy is absorbed by retinal tissue and converted to heat, coagulating tissue and blood vessels. This seals leaking vessels, closes abnormal new vessels, and creates chorioretinal adhesions around tears that prevent retinal detachment.
\section*{History}
Retinal laser photocoagulation was developed by Franciscus Donders' successors and others from the 1960s, and large trials in the 1970s established it as standard treatment for proliferative diabetic retinopathy and macular edema.
\section*{Why This Test or Procedure Is Ordered}
Ordered for proliferative diabetic retinopathy, clinically significant diabetic macular edema, retinal tears, central serous retinopathy, and certain vascular retinal diseases to prevent or limit vision loss.
\section*{Is This a Primary or Secondary Test or Procedure}
A primary therapeutic procedure for many retinal diseases, used alone or with other treatments such as anti-VEGF injections.
\section*{How This Test or Procedure Is Performed}
After the pupil is dilated and eye drops applied, a contact lens is placed on the eye and laser pulses are delivered through a slit lamp to the target areas of the retina. The number and pattern of burns depend on the condition being treated.
\section*{What Preparation Is Needed}
Dilation of the pupil with eye drops and review of the patient's medications, including blood thinners, which may be continued in most cases.
\section*{Contraindications and Precautions}
Precautions include media opacity such as cataract or vitreous hemorrhage that blocks the laser beam, and caution with perifoveal treatment, which risks central scotoma.
\section*{Risks}
Temporary blurred vision, discomfort, scarring, decreased night or peripheral vision after panretinal photocoagulation, and rarely bleeding, retinal detachment, or increased eye pressure.
\section*{What Happens After the Test or Procedure}
The patient may experience transient visual blurring and uses prescribed eye drops. Vision may take weeks to stabilize, and follow-up examinations monitor the response.
\section*{Understanding Your Results}
Control of the treated condition, with reduced leakage and regression of new vessels, indicates success. Persistent or progressive disease may require additional laser or other treatment.
\section*{Normal Results}
Stabilized or improved vision with sealing of leaking vessels, regression of abnormal new vessels, and no evidence of ongoing retinal disease progression.
\section*{Follow-up Tests and Procedures}
Regular dilated eye examinations are required, and further laser sessions, injections, or vitrectomy may be offered depending on the response.
\section*{References}
\begin{enumerate}
  \item MedlinePlus. Laser eye surgery for diabetic retinopathy. U.S. National Library of Medicine; 2025.
  \item Current Medical Diagnosis and Treatment. McGraw-Hill; 2025.
\end{enumerate}

\clearpage
